\section{Related Work}
\label{sec:related}

Having established BAA measurement as an empirically open question, we review three bodies of literature that each address part of the problem --- and show why none closes the gap this paper targets.

\subsection{AI-to-Human Alignment: Well-Measured, One-Directional}

The dominant paradigm in AI alignment research operationalizes ``alignment'' as the degree to which AI outputs match human preferences, values, or intent. RLHF \citep{Ouyang2022Training, Bai2022Training} provides preference-labeled training data collected from human raters. Benchmarks such as HELM \citep{Liang2022Holistic}, TruthfulQA \citep{Lin2022TruthfulQA}, and BIG-Bench \citep{Srivastava2022Beyond} measure AI performance on tasks that proxy human-valued behaviors. This literature is mature, with standardized evaluation infrastructure and reproducible scores across model generations.

A critical feature of this paradigm is that it treats human evaluators as a fixed, stable signal source. The implicit assumption is that rater behavior is stationary: raters in 2024 apply the same judgment standards as raters in 2022. Whether this assumption holds as raters accumulate experience with AI-generated text has not been systematically examined.

\subsection{Human Behavioral Adaptation: Qualitative Evidence, No Measurement Framework}

A smaller literature documents how human behavior adapts in response to AI quality. \citet{Dell2023Navigating} study management consultants using AI assistance and find that AI-assisted workers progressively delegate cognitively demanding tasks to AI, producing measurable deskilling in unassisted task performance. Their effect sizes (0.3--0.5 SD) motivate our choice of $|\tau| \geq 0.2$ as a minimum detectable effect threshold. Crucially, \citet{Dell2023Navigating} measure task assignment behavior in a controlled experimental setting, not behavioral signals in naturalistic interaction logs.

\citet{Perez2023Sycophancy} demonstrate that AI sycophancy suppresses human expression of disagreement --- complementary to BAA: if AI models are trained to suppress explicit corrections, measuring correction frequency as a BAA proxy will undercount genuine disengagement.

\citet{Shen2024Towards} provide the most direct conceptual ancestor of the present work, surveying bidirectional human-AI alignment and identifying the measurement gap at the framework level. They propose behavioral proxy dimensions but provide no empirical measurement on public interaction logs. We implement and test their implied measurement program, producing --- to our knowledge --- the first large-scale empirical test.

\subsection{Temporal Analysis of AI Interaction Logs}

The Chatbot Arena platform \citep{Zheng2023Judging} provides timestamped human preference votes across model pairs, making it the natural data source for temporal preference entropy analysis. We were unable to access the primary \texttt{lmsys/chatbot\_arena\_conversations} dataset due to gating; the publicly available fallback contains no timestamps.

WildChat \citep{Zhao2024WildChat} provides 1 million ChatGPT conversation logs with IP-hash pseudonymization, timestamps, and topic tags. Prior analyses focus on content characterization --- not on behavioral trend analysis within returning-user cohorts. We introduce this analytical lens.

Mann-Kendall trend analysis \citep{Mann1945Nonparametric, Kendall1975Rank} is standard in environmental and behavioral time series analysis. The Hamed-Rao modification \citep{Hamed1998Modified} corrects for autocorrelated series; it is critical in our context (ACF lag-1 = 0.634). To our knowledge, no prior work has applied Hamed-Rao corrected Mann-Kendall to behavioral proxies in AI interaction logs.

\subsection{Our Position}

Existing work on AI-to-human alignment does not account for human behavioral adaptation over time. Existing work on human behavioral adaptation uses controlled experiments, not public logs. Existing temporal analyses of WildChat and LMSYS Arena do not construct returning-user cohorts or test behavioral trends. \citet{Shen2024Towards} identify the BAA framework conceptually but provide no empirical measurement. We occupy the gap.
